\section{Introducere}

\emph{Eye-tracking}-ul permite observarea directă a modului în care o persoană explorează vizual o scenă. În loc să analizăm doar rezultatul final al unei sarcini, putem urmări traseul privirii, zonele în care participantul revine, timpul până la prima fixație pe țintă și felul în care atenția se mută între regiuni ale imaginii \cite{yarbus1967,just1980}. În cazul jocului \emph{Find Waldo}, aceste informații sunt utile deoarece participanții trebuie să caute activ un personaj într-o imagine aglomerată. Astfel, putem analiza atât comportamente orientate de scopul sarcinii, de tip \emph{top-down}, cât și comportamente influențate de elemente vizuale evidente, de tip \emph{bottom-up} \cite{itti1998,otero2016waldo,wolfe2020visual}. Ulterior, traseele vizuale extrase pot fi transformate în serii temporale și în date structurate, pe baza cărora se pot aplica metode de analiză și modelare.

Prezenta lucrare de licență se încadrează în domeniile analizei de date și a învățării automate (\emph{Machine Learning}). Alegerea acestei teme a venit ca o urmare a participării de la Noaptea Cercetătorilor și din necesitatea de a înțelege datele colectate. Din acest motiv, contribuția proprie a început cu mult înaintea scrierii codului, prin pregătirile realizate pentru eveniment, unde am testat mai multe experimente. Inițial, am luat în considerare sarcini de citire și înțelegere a textului, apoi urmărirea unor instrucțiuni audio pe materiale video sau teste de tip \emph{human benchmark}. În cele din urmă, am ales experimentul \emph{Find Waldo}, deoarece se potrivea cel mai bine contextului: este o activitate interactivă pentru un public divers (copii și adulți), oferind în același timp posibilitatea de a colecta date relevante pentru studiu. 

Rolul meu a fost unul tehnic, concentrat pe configurarea modului de lucru, calibrarea echipamentului și colectarea efectivă a datelor. În plus, pentru public am creat un \emph{setup} interactiv cu multiple ecrane. Unul dintre ecrane afișa perspectiva participantului care purta \emph{eye-tracker}-ul, având suprapus un indicator ce reprezenta locul în care privește acesta. Un al doilea ecran rula un \emph{dashboard} cu grafice actualizate în timp real, ce oferea statistici precum: coordonatele privirii, viteza de mișcare oculară, fixațiile și sacadele. De asemenea, la finalul unor sesiuni, prezentam participanților care erau curioși de rezultate un feedback direct folosind \emph{Neon Player}, înlocuind temporar graficele cu hărți de atenție (\emph{heatmaps}) și traiectoriile vizuale (\emph{gazepaths}) generate pe baza nivelurilor tocmai încheiate. Astfel, munca mea a asigurat atât infrastructura pentru afișarea informațiilor către public, cât și colectarea datelor. 

\subsection{Obiectivele tezei}

Teza urmărește două obiective principale. Primul este construirea unui sistem de prelucrare pentru datele colectate cu \emph{eye-tracker}-ul, pornind de la înregistrări și ajungând la un set de date segmentat pe niveluri, verificat manual și descris prin metrici numerice. Acest pas este important deoarece datele nu provin dintr-un experiment controlat de laborator, ci dintr-un eveniment public, unde apar inevitabil sesiuni incomplete, zgomot, diferențe între participanți și momente în care calitatea înregistrării scade.

Al doilea obiectiv este folosirea acestor metrici pentru a analiza comportamentul vizual. Ne interesează, pe de o parte, dacă participanții urmăresc zonele evidente ale imaginii sau dacă își controlează privirea în funcție de scopul sarcinii. Pe de altă parte, verificăm dacă primele secunde ale unei sesiuni conțin suficientă informație pentru a anticipa succesul, durata căutării sau următoarea zonă privită.

Structura tezei urmărește pașii prin care am ajuns de la colectarea datelor la analiza și modelarea lor. \emph{Secțiunea 2} introduce contextul și conceptele teoretice legate de atenția vizuală și mișcările oculare. \emph{Secțiunea 3} prezintă detalii despre organizarea experimentului și procesul de colectare a datelor. \emph{Secțiunea 4} descrie pașii de adnotare și verificare a înregistrărilor, pas necesar pentru validarea și analiza datelor. După obținerea datelor curate, \emph{Secțiunea 5}, respectiv \emph{Secțiunea 6} cuprind etapele de preprocesare, analiza exploratorie a datelor (\emph{EDA}) și etapele ulterioare de analiză. În aceste secțiuni, voi detalia metricile extrase, comportamentul și tiparele observate, voi face comparația bazată pe \emph{saliency} și voi discuta concluziile rezultate din teste. Secțiunea 7 analizează strategiile de căutare, investigând dacă participanții își păstrează un stil individual de explorare vizuală, independent de dificultatea inerentă a nivelului. \emph{Secțiunea 8} prezintă în detaliu modelele de inteligență artificială folosite pe datele extrase pentru a rezolva task-urile specifice, comparații între modurile lor de funcționare și rezultatele pe datele noastre. Secțiunea 9 prezintă analiza de interpretabilitate și stabilitate, prin care verificăm ce caracteristici influențează deciziile modelelor și cât de stabile sunt rezultatele la eliminarea datelor zgomotoase. În final, Secțiunea 10 sintetizează concluziile și limitările întregului proiect.


\subsection{Contribuții}

Contribuțiile sunt împărțite în trei direcții:

Pe partea experimentală, m-am implicat direct în dezvoltarea modului de colectare a datelor la Noaptea Cercetătorilor. După eveniment, am construit pașii de segmentare, adnotare și verificare manuală, pas esențial prin care am reușit să transform înregistrările obținute de la \emph{eye-tracker} într-un set de date utilizabil.

Din punct de vedere al analizei, am realizat explorarea datelor prin extragerea unor metrici spațio-temporale: timpul până la prima fixație pe Waldo, entropia spațială, entropia de tranziție, acoperirea spațială și raportul fixațiilor pe Waldo. Am folosit aceste metrici pentru a analiza de ce anumite scene din joc sunt mai dificile și modul cum își schimbă participanții strategiile de explorare în funcție de nivel. În aceeași etapă, am generat hărți de atenție și le-am comparat cu hărțile de \emph{saliency} pentru a evalua raportul dintre reflexele vizuale și controlul cognitiv orientat pe scopul de a găsi personajul. 

Mai departe, m-am ocupat de partea de modelare a datelor. În acest sens, am construit mai multe modele de inteligență artificială: tabulare (folosind Regresia Logistică și rețelele MLP), recurente (\emph{BiLSTM} pe ferestre de timp) și modele Markov de ordin întâi pentru anticiparea următoarei regiuni la care se va uita un participant. Pentru fiecare variantă am rulat experimente și am urmărit dacă rezultatele rămân stabile. Pe lângă acestea, am evaluat diferențele stilurilor de explorare vizuală specifice fiecărui participant. Mai mult, am adăugat o analiză de interpretabilitate (\emph{feature and sensitivity analysis}) pentru modelele antrenate, care arată ce atribute influențează deciziile algoritmilor. Astfel, rezultatele obținute pot fi interpretate și din perspectiva comportamentului vizual, nu doar ca scoruri de clasificare.